\section{Los tres grandes de China y la ventaja en costos}

En este episodio también se menciona a los tres grandes de China. Aunque la transcripción no entra en muchos detalles, por el contexto anterior y posterior la ventaja de las empresas chinas de modelos parece venir más bien de los costos, la velocidad de ingeniería, los escenarios de aplicación y el ecosistema de desarrolladores, y no de aplastar a la competencia en un solo benchmark. Los escenarios de Coding/Agent son especialmente sensibles al costo, porque las tareas de largo alcance multiplican cada llamada al modelo.

\subsection{Por qué el costo es una variable estratégica}

En un escenario de chat, quizá el usuario todavía acepte un modelo algo más caro; en un escenario de Agent, una sola tarea puede requerir decenas de rondas de llamadas, ejecución de herramientas y verificación. Si el costo por llamada es demasiado alto, al producto le cuesta escalar. Si los modelos nacionales baratos y eficaces logran conectarse a un buen harness, podrían obtener una ventaja en un gran volumen de tareas de frecuencia media y alta.

\begin{importantbox}{La ventaja en costos no es una ventaja de gama baja}
Un modelo barato no está limitado a tareas de poco valor. Un framework de Agent puede usar el modelo barato en los pasos de alta frecuencia y el modelo potente en las decisiones críticas. El enrutamiento entre varios modelos convierte la estructura de costos en un factor de competitividad.
\end{importantbox}

\subsection{Resumen de la sección}

La oportunidad de las empresas chinas de modelos consiste en lo siguiente: si logran incorporar su ventaja en costos, su velocidad de ingeniería y la retroalimentación de sus aplicaciones a los flujos de trabajo de Coding/Agent, podrían ocupar una posición diferenciada en el segundo acto.
